\documentclass[10pt,a4paper]{amsart}
\usepackage[margin=19mm]{geometry}
\usepackage{amsmath,amssymb,mathtools}
\usepackage[scr=rsfs,cal=euler]{mathalfa}
\usepackage{xfrac}
\usepackage[unicode,hidelinks,bookmarks=false]{hyperref}
\newcommand{\ie}{i.e.\ }
\newcommand{\cf}{cf.\ }
\newcommand{\resp}{respectively\ }
\newcommand{\Subsection}{\subsection}
\providecommand{\nobreakdash}{\nobreak\hbox{-}}
\setlength{\emergencystretch}{2em}
\multlinegap0pt
\usepackage{amssymb}
\usepackage{xcolor}
\newtheorem{proposition}{Proposition}
\newcommand{\RR}{\mathbb{R}}
\newcommand{\ZZ}{\mathbb{Z}}
\newcommand{\calD}{\mathcal{D}}
\newcommand{\calO}{\mathcal{O}}
\title{Pseudo-Differential Operators and Generalized Random Fields over Tori}
\author{Nicolas Escobar-Velasquez}
\date{Source: arXiv:2511.09423v2 (CC BY 4.0)}
\begin{document}
\maketitle

\noindent\emph{Source and licence.} Pseudo-Differential Operators and Generalized Random Fields over Tori. Version arXiv:2511.09423v2 (CC BY 4.0), distributed by arXiv under the Creative Commons Attribution 4.0 International licence, http://creativecommons.org/licenses/by/4.0/, as recorded in the arXiv licence metadata for this exact version and shown on the abstract page https://arxiv.org/abs/2511.09423v2. Full licence notice: \texttt{licenses/arxiv-2511.09423-CC-BY-4.0.txt}.\par\smallskip

\noindent\textit{Editorial scope.} Counted unit: the article's proposition prop:canonical-cov (source0.tex lines 829--841). The article's complete original proof of it is delivered with it (source0.tex lines 843--873, one \texttt{proof} environment of 31 lines). No mathematics is written, repaired or re-derived: the statement and the proof are the article's own lines, in the article's own notation, and no retained line is substituted or re-flowed.  Retained uncounted as the source-local context the proof needs: lines 802--804; lines 811--813.  Nothing was omitted from any retained span, so the package carries no omission record.  No retained line contains a citation, so the wrapper carries no bibliography and no bibliography entry.  No retained line contains a figure environment, an image-inclusion command or drawing code, so no image is delivered and the package declares no asset dependency.  Editorial formatting only, in addition: an AMS \texttt{amsart} preamble that restates the article's own declarations and macros used by the retained lines, namely the environments \texttt{proposition} on separate counters, together with the standard packages those lines need.  The retained environments print in this document as Proposition 1 in order of appearance, whereas the article prints the same items with the numbering of its own full document.  The complete native source of the article as distributed in the arXiv e-print for this exact version is bundled unchanged as \texttt{sources/188981/source0.tex}.
\begin{equation}\label{eq:green}
    -\Delta\, G_{-\Delta} \;=\; \delta_0 - 1 ,
\end{equation}

    \begin{equation}\label{eq:canonical-green}
        C_G \;=\; 4\pi^2\, G_{-\Delta} .
    \end{equation}

\begin{proposition}[Covariance of the canonical field]\label{prop:canonical-cov}
The covariance distribution of $G$ is
$C_G(h) = \sum_{k \in \ZZ^d\setminus\{0\}} |k|^{-2} e^{2\pi i k\cdot h}$, and
\begin{enumerate}
\item[(i)] for $d=1$ and $h \in [-\tfrac12,\tfrac12]$,
    $\;C_G(h) = \tfrac{\pi^2}{3} - 2\pi^2|h| + 2\pi^2 h^2$;
\item[(ii)] for $d=2$, $\;C_G(h) = -2\pi\log|h| + O(1)$ as $h \to 0$;
\item[(iii)] for $d \geq 3$,
    $\;C_G(h) = c_d|h|^{2-d} + o\bigl(|h|^{2-d}\bigr)$ as $h \to 0$, with
    $c_d = \pi^{2-d/2}\,\Gamma\bigl(\tfrac d2 - 1\bigr)$.
\end{enumerate}
In particular $C_G$ is a bounded continuous function if and only if $d = 1$.
\end{proposition}

\begin{proof}
Part (i) is the Fourier series of the second Bernoulli polynomial,
\[
    \sum_{k\geq1} k^{-2}\cos(2\pi k h) \;=\; \pi^2\bigl(h^2-h+\tfrac16\bigr),
    \qquad h \in [0,1],
\]
of which $C_G(h)$ is twice the sum.

For (ii) and (iii) we compare \eqref{eq:green} with its Euclidean counterpart.
Let $E$ be the Newtonian kernel on $\RR^d$, the fundamental solution of
$-\Delta$:
\[
    E(h) \;=\; -\tfrac{1}{2\pi}\log\|h\| \quad (d=2), \qquad
    E(h) \;=\; \frac{\Gamma\bigl(\tfrac d2-1\bigr)}{4\pi^{d/2}}\,\|h\|^{2-d}
    \quad (d\geq3) .
\]
The torus is flat, so a ball $B$ about the origin of radius $<\tfrac12$ is
isometric to a Euclidean ball and $-\Delta(4\pi^2 E) = 4\pi^2\delta_0$ holds in
$\calD'(B)$. Put $Q(h) := \tfrac{2\pi^2}{d}\|h\|^2$, so that
$-\Delta Q = -4\pi^2$. By \eqref{eq:green} and \eqref{eq:canonical-green},
\[
    -\Delta\bigl(C_G - 4\pi^2E - Q\bigr) \;=\; 4\pi^2\delta_0 - 4\pi^2
        - 4\pi^2\delta_0 + 4\pi^2 \;=\; 0 \qquad\text{in } \calD'(B) ,
\]
so $C_G - 4\pi^2E - Q$ is a harmonic distribution on $B$ and hence, by Weyl's
lemma, a real-analytic function there; in particular it is bounded near $0$.
Therefore $C_G(h) = 4\pi^2 E(h) + \calO(1)$ as $h \to 0$, and multiplying the
displayed constants by $4\pi^2$ gives (ii) and (iii), the error being $\calO(1)$
and hence $o(|h|^{2-d})$ when $d\geq3$. Boundedness fails for $d \geq 2$ because
$C_G(0) = \sum_{k \neq 0}|k|^{-2}$ diverges there.
\end{proof}

\end{document}
